\chapter{Experimentación y evaluación}

\section{Objetivos experimentales}
\pending{Definir qué preguntas responde cada experimento y distinguir validación técnica de evaluación científica.}

\section{Protocolo experimental}
\pending{Documentar datasets, particiones, semillas, codificaciones, entrenamiento, métricas y condiciones de ejecución.}

\section{Validación técnica de la QCNN}
\pending{Presentar pruebas de codificación, circuitos, pooling, gradientes, entrenamiento, shots, ruido y reproducibilidad.}

\section{Experimentos con Iris}
\pending{Presentar la convergencia exacta y el diagnóstico inicial sobre Iris binario.}

\section{Experimentos con Breast Cancer Wisconsin}
\pending{Describir el dataset mayor, la partición 70/15/15 y su finalidad.}

\section{Preprocesamiento y reducción mediante PCA}
\pending{Explicar la reducción de 30 características a cuatro componentes y la prevención de fuga de información.}

\section{Convergencia}
\pending{Analizar 10, 20 y 30 épocas y justificar la configuración de referencia.}

\section{Variabilidad entre semillas}
\pending{Analizar la estabilidad de la QCNN en varias inicializaciones.}

\section{Comparación de codificaciones}
\pending{Comparar codificación angular y por fase manteniendo constantes los demás factores.}

\section{Shots y ruido}
\pending{Analizar el efecto de 50 shots y ruido depolarizante sobre rendimiento, estabilidad y coste.}

\section{Comparación con el baseline clásico}
\pending{Comparar con regresión logística usando la misma partición y métricas.}

\section{Evaluación de la plataforma}
\pending{Evaluar usabilidad, trazabilidad, facilidad de repetición y extensibilidad cuando la aplicación esté integrada.}

\section{Amenazas a la validez}
\pending{Analizar tamaño muestral, número de semillas, PCA, simulación, hardware, partición y coste computacional.}
